\chapter{Menstruation}

\section*{Definition and Overview}
Menstruation, commonly referred to as a period, is the cyclical shedding of the functional layer of the endometrium (uterine lining) accompanied by bleeding, which occurs in healthy females of reproductive age who are not pregnant. This physiological process is regulated by complex interactions within the hypothalamic-pituitary-ovarian (HPO) axis. Menstruation marks the start of the menstrual cycle, which typically averages 28 days but can normally range from 21 to 35 days. The onset of menstruation (menarche) usually occurs between ages 12 and 13, and its permanent cessation (menopause) occurs on average around age 51. Understanding menstruation is fundamental to reproductive health, family planning, and the diagnosis of gynaecological disorders.

\section*{Epidemiology}
Menstruation affects approximately 50\% of the global population for a significant portion of their lives, spanning roughly 35 to 40 years. In Australia, millions of women, girls, and gender-diverse individuals menstruate. Normal menstrual bleeding lasts between 3 and 7 days, with an average blood loss of 30 to 40 millilitres per cycle. Disorders of menstruation are highly prevalent; abnormal uterine bleeding (AUB) affects up to 30\% of women of reproductive age globally, and dysmenorrhea (painful periods) is reported by up to 70\% to 90\% of young women, representing a major cause of school and work absenteeism.

\section*{Aetiology and Pathophysiology}
Menstruation occurs as a result of cyclical hormonal changes in the absence of conception and implantation. The physiological steps include:
\begin{itemize}
    \item \textbf{Follicular phase and ovulation:} Follicle-stimulating hormone (FSH) stimulates ovarian follicle development, which secretes oestrogen to proliferate the endometrium. Ovulation occurs mid-cycle triggered by a luteinising hormone (LH) surge.
    \item \textbf{Luteal phase and progesterone secretion:} The ruptured follicle becomes the corpus luteum, secreting progesterone to prepare the endometrium for implantation.
    \item \textbf{Hormone withdrawal:} If fertilisation does not occur, the corpus luteum degenerates, leading to a sudden drop in progesterone and oestrogen levels.
    \item \textbf{Endometrial shedding:} Progesterone withdrawal triggers local vasoconstriction of spiral arteries, enzymatic breakdown of the endometrial matrix, tissue necrosis, and shedding of the stratum functionalis layer.
    \item \textbf{Myometrial contraction:} Prostaglandins released from the shedding tissue cause uterine contractions to expel the menstrual blood and limit blood loss, which also mediates normal menstrual cramping.
\end{itemize}

\section*{Clinical Features}
A normal menstrual period presents with predictable features, but can be accompanied by physical and emotional symptoms:
\begin{itemize}
    \item \textbf{Menstrual bleeding:} Discharge of blood, tissue debris, and mucus from the vagina, varying from light spotting to heavier flow, without large blood clots.
    \item \textbf{Pelvic cramping (dysmenorrhea):} Mild to moderate aching or cramping in the lower abdomen, pelvis, or lower back.
    \item \textbf{Premenstrual symptoms (PMS):} Bloating, breast tenderness, mild weight gain, acne breakouts, fatigue, and mood fluctuations (irritability, anxiety) in the days preceding flow.
    \item \textbf{Changes in vaginal discharge:} Transition from thick, sticky mucus to watery, clear mucus during ovulation, and subsequent post-ovulatory changes.
    \item \textbf{Normal cyclical parameters:} Cycle length of 21--35 days, bleeding duration of 3--7 days, and predictable recurrence.
\end{itemize}

\section*{Differential Diagnosis}
Deviations from normal menstruation (abnormal uterine bleeding or severe pain) require differentiation from several pathological conditions:
\begin{itemize}
    \item \textbf{Pregnancy and ectopic pregnancy:} Must always be ruled out in cases of delayed, missed, or abnormal bleeding.
    \item \textbf{Uterine fibroids (leiomyomas) or polyps:} Benign growths causing heavy menstrual bleeding (menorrhagia) or intermenstrual bleeding.
    \item \textbf{Endometriosis or adenomyosis:} Pathologies presenting with progressive, severe dysmenorrhea, chronic pelvic pain, and deep dyspareunia.
    \item \textbf{Polycystic ovary syndrome (PCOS):} An endocrine disorder presenting with oligomenorrhea (infrequent periods), amenorrhea, and hyperandrogenism.
    \item \textbf{Thyroid disease:} Hypothyroidism or hyperthyroidism can disrupt the HPO axis, causing irregular, heavy, or absent menstruation.
    \item \textbf{Coagulation disorders:} Von Willebrand disease or thrombocytopenia presenting with heavy menstrual bleeding since menarche.
\end{itemize}

\section*{When to Seek Medical Care}
Menstruating individuals should seek medical advice if they experience significant changes in their normal cycle, severe pain, or bleeding that interferes with daily activities.

\textbf{Call 000 (or local emergency services) for:}
\begin{itemize}
    \item \textbf{Severe, acute lower abdominal pain:} Especially if accompanied by shoulder tip pain, dizziness, or fainting, which can indicate a ruptured ectopic pregnancy.
    \item \textbf{Signs of toxic shock syndrome (TSS):} Sudden high fever, vomiting, watery diarrhoea, sunburn-like rash, and low blood pressure, associated with tampon use.
    \item \textbf{Severe haemorrhage:} Bleeding so heavy that it causes dizziness, confusion, or shortness of breath.
\end{itemize}
Other reasons to see a doctor include periods lasting longer than 8 days, cycles shorter than 21 days or longer than 35 days, bleeding between periods or after sexual intercourse, or absence of periods for more than 3 months.

\section*{Diagnosis and Investigations}
Investigations are directed at identifying the cause of abnormal or painful menstruation:
\begin{itemize}
    \item \textbf{Pregnancy test (Urine or serum hCG):} The initial screening test for any abnormal bleeding or missed period.
    \item \textbf{Pelvic ultrasound:} The primary imaging test to evaluate uterine structure, identify fibroids, polyps, or adenomyosis, and assess the ovaries.
    \item \textbf{Blood tests:} Full blood count (to check for anaemia), iron studies, thyroid-stimulating hormone (TSH), prolactin, and coagulation profiles.
    \item \textbf{Hormonal panel:} FSH, LH, oestrogen, progesterone, and free testosterone to evaluate ovulation or suspected PCOS.
    \item \textbf{Endometrial biopsy:} Performed in women over 35 or those with risk factors to rule out endometrial hyperplasia or malignancy.
\end{itemize}

\section*{Management}
Normal menstruation requires hygienic care, while menstrual disorders are managed pharmacologically or surgically:
\begin{itemize}
    \item \textbf{Menstrual hygiene products:} Safe use of sanitary pads, tampons, menstrual cups, or period underwear. Tampons must be changed every 4 to 8 hours to prevent TSS.
    \item \textbf{Pharmacological management of pain:} NSAIDs (e.g. ibuprofen, mefenamic acid) block prostaglandin synthesis and are highly effective for dysmenorrhea.
    \item \textbf{Hormonal therapies:} Combined oral contraceptive pills, progestogen-only pills, contraceptive implants, or the levonorgestrel-releasing intrauterine device (IUD) to regulate cycles, reduce bleeding volume, and alleviate pain.
    \item \textbf{Surgical options:} Hysteroscopic resection of polyps or fibroids, endometrial ablation, or hysterectomy for severe, refractory cases of AUB or adenomyosis.
\end{itemize}

\section*{Complications}
Complications related to menstrual disorders include:
\begin{itemize}
    \item \textbf{Iron-deficiency anaemia:} Secondary to chronic heavy menstrual blood loss, presenting with fatigue, pallor, and weakness.
    \item \textbf{Infertility:} Associated with underlying causes of absent or irregular menstruation (e.g. anovulation, PCOS, endometriosis).
    \item \textbf{Osteopenia and osteoporosis:} Long-term complication of prolonged amenorrhea (lack of oestrogen), leading to reduced bone density.
    \item \textbf{Toxic shock syndrome (TSS):} A rare, life-threatening bacterial infection associated with prolonged tampon use.
    \item \textbf{Reduced quality of life:} Chronic pain and heavy bleeding impacting social, educational, and occupational activities.
\end{itemize}

\section*{Prognosis and Outcomes}
The prognosis for managing menstrual disorders is generally excellent. The majority of menstrual irregularities and painful periods respond well to conservative lifestyle adjustments, targeted hormonal therapies, or minor surgical interventions. For most individuals, menstrual cycles become more regular in their 20s and 30s before becoming irregular again during the perimenopausal transition. Education, open discussions, and early medical intervention can effectively prevent long-term complications like severe anaemia or bone loss, promoting excellent reproductive health.

\section*{Prevention and Self-Care}
While menstruation itself is a natural process, discomfort and menstrual disorders can be managed or prevented with self-care:
\begin{itemize}
    \item \textbf{Track your cycle:} Use a calendar or app to monitor cycle lengths, bleeding intensity, and pain levels, which assists in early detection of abnormalities.
    \item \textbf{Apply local heat:} Use a heat pack or hot water bottle on the lower abdomen to relieve uterine cramping.
    \item \textbf{Regular moderate exercise:} Physical activity can reduce pain and improve mood by releasing endorphins and lowering prostaglandin levels.
    \item \textbf{Maintain a balanced diet:} Ensure adequate intake of iron-rich foods (e.g. lean meats, leafy green vegetables) to offset menstrual blood loss.
    \item \textbf{Practice safe hygiene:} Change tampons and sanitary pads regularly, wash the external genital area with water only, and avoid douching.
\end{itemize}

\section*{Sources and Further Reading}
The following reputable organisations provide further information on menstruation:
\begin{itemize}
    \item Jean Hailes for Women's Health. \textit{Menstrual cycle, periods, and management of disorders.} https://www.jeanhailes.org.au/
    \item Better Health Channel. \textit{Menstrual cycle: what is normal and when to get help.} https://www.betterhealth.vic.gov.au/
    \item Healthdirect Australia. \textit{Menstruation and period problems overview.} https://www.healthdirect.gov.au/
    \item Mayo Clinic. \textit{Menstrual cycle: what's normal, what's not.} https://www.mayoclinic.org/
    \item NHS. \textit{Periods: advice, sanitary products, and abnormalities.} https://www.nhs.uk/
\end{itemize>
